
\documentclass[11pt,a4paper]{article}

\usepackage[T2A]{fontenc}
\usepackage[utf8]{inputenc}
\usepackage[english,russian]{babel}
\usepackage{amsmath,amssymb,mathtools}
\usepackage{booktabs,array}
\usepackage{geometry}
\usepackage{enumitem}
\usepackage[hidelinks]{hyperref}

\geometry{margin=2cm}
\setlength{\parindent}{1.25em}
\setlength{\parskip}{0.35em}

\newcommand{\R}{\mathbb{R}}
\newcommand{\norm}[1]{\left\lVert #1\right\rVert_2}
\newcommand{\ip}[2]{\left\langle #1,#2\right\rangle}
\newcommand{\argmin}{\operatorname*{arg\,min}}
\newcommand{\argmax}{\operatorname*{arg\,max}}
\newcommand{\tr}{\operatorname{tr}}
\newcommand{\diag}{\operatorname{diag}}
\newcommand{\matvec}{\operatorname{matvec}}
\newcommand{\sign}{\operatorname{sign}}

\title{Выбор параметра ridge-регуляризации\\
в matrix-free single-index солвере\\
методом Hybrid Golub--Kahan и projected GCV}
\author{}
\date{}

\begin{document}
\maketitle

\section{Постановка задачи}

На фиксированном шаге чередующейся оптимизации заданы локальные статистики
\[
    I_j\in\R^{P},
    \qquad
    U_j\in\R^{P\times d},
    \qquad
    j=1,\ldots,J,
\]
и локальные коэффициенты \(c_j\in\R\).
Глобальное направление \(\beta\in\R^d\) вычисляется из ridge-задачи
\begin{equation}
\label{eq:ridge}
    \widehat\beta_\lambda
    =
    \argmin_{\beta\in\R^d}
    \left\{
        \sum_{j=1}^{J}
        \norm{I_j-c_jU_j\beta}^{2}
        +
        \lambda\norm{\beta-\beta_0}^{2}
    \right\},
\end{equation}
где \(\beta_0\) --- направление с предыдущего внутреннего или внешнего шага.

Определим
\[
    A=
    \begin{pmatrix}
        c_1U_1\\
        \vdots\\
        c_JU_J
    \end{pmatrix}
    \in\R^{JP\times d},
    \qquad
    b=
    \begin{pmatrix}
        I_1\\
        \vdots\\
        I_J
    \end{pmatrix}
    \in\R^{JP}.
\]
Тогда \eqref{eq:ridge} принимает вид
\begin{equation}
\label{eq:ridge-compact}
    \widehat\beta_\lambda
    =
    \argmin_{\beta}
    \left\{
        \norm{A\beta-b}^{2}
        +
        \lambda\norm{\beta-\beta_0}^{2}
    \right\}.
\end{equation}

Матрица \(A\) не формируется. Достаточно операторов
\begin{align}
\label{eq:matvec}
    Av
    &=
    \begin{pmatrix}
        c_1U_1v\\
        \vdots\\
        c_JU_Jv
    \end{pmatrix},
    \\
\label{eq:rmatvec}
    A^\top z
    &=
    \sum_{j=1}^{J}
    c_jU_j^\top z_j,
    \qquad
    z=
    (z_1^\top,\ldots,z_J^\top)^\top .
\end{align}
Одна пара операций \(Av\) и \(A^\top z\) требует
\[
    O(JPd)
\]
арифметических операций и сохраняет разреженную или потоковую реализацию локальных статистик.

Если текущая модель содержит локальный свободный коэффициент \(a_jS_j\), следует заменить
\[
    I_j
    \quad\text{на}\quad
    I_j-a_jS_j.
\]
Остальная схема не меняется.

\section{Соответствие между \(\lambda\) и \texttt{damp}}

LSMR с параметром \texttt{damp}\(=\mu\) решает
\begin{equation}
\label{eq:lsmr-damp}
    \min_x
    \left\{
        \norm{Ax-b}^{2}
        +
        \mu^2\norm{x}^{2}
    \right\}.
\end{equation}
Поэтому параметр \(\lambda\) из \eqref{eq:ridge-compact} нельзя напрямую передавать как
\texttt{damp}.

Для центрированного штрафа вводится
\[
    \delta=\beta-\beta_0,
    \qquad
    r_0=b-A\beta_0.
\]
Тогда
\begin{equation}
\label{eq:shifted}
    \widehat\delta_\lambda
    =
    \argmin_{\delta}
    \left\{
        \norm{A\delta-r_0}^{2}
        +
        \lambda\norm{\delta}^{2}
    \right\},
    \qquad
    \widehat\beta_\lambda=\beta_0+\widehat\delta_\lambda.
\end{equation}
Для обычного вызова LSMR требуется
\begin{equation}
\boxed{
    \texttt{damp}=\sqrt{\lambda}.
}
\end{equation}

Эквивалентная расширенная система имеет вид
\[
    \min_{\beta}
    \left\|
    \begin{pmatrix}
        A\\
        \sqrt{\lambda}I_d
    \end{pmatrix}\beta
    -
    \begin{pmatrix}
        b\\
        \sqrt{\lambda}\beta_0
    \end{pmatrix}
    \right\|_2^2.
\]
В этом варианте регуляризация уже включена в оператор, поэтому
\begin{equation}
\boxed{
    \texttt{damp}=0.
}
\end{equation}
Одновременное использование расширенной системы и ненулевого \texttt{damp}
добавляет второй штраф относительно нуля.

\section{Hybrid Golub--Kahan regularization}

\subsection{Одна Krylov-проекция для всех значений \(\lambda\)}

Для задачи \eqref{eq:shifted} выполняется \(k\) шагов билинеаризации Голуба--Кахана:
\begin{equation}
\label{eq:gkb}
    AV_k=U_{k+1}B_k,
    \qquad
    r_0=\rho u_1,
    \qquad
    \rho=\norm{r_0},
\end{equation}
где
\[
    V_k\in\R^{d\times k},
    \qquad
    U_{k+1}\in\R^{JP\times(k+1)},
    \qquad
    B_k\in\R^{(k+1)\times k}
\]
и \(B_k\) --- нижняя бидиагональная матрица.

При ограничении \(\delta=V_ky\) исходная задача заменяется малой ridge-задачей
\begin{equation}
\label{eq:projected}
    y_{\lambda,k}
    =
    \argmin_{y\in\R^k}
    \left\{
        \norm{B_ky-\rho e_1}^{2}
        +
        \lambda\norm{y}^{2}
    \right\}.
\end{equation}
После этого
\begin{equation}
\label{eq:beta-projected}
    \widetilde\beta_{\lambda,k}
    =
    \beta_0+V_ky_{\lambda,k}.
\end{equation}

Один Krylov-процесс обслуживает всю сетку значений \(\lambda\). Повторные полные
запуски LSMR для каждого кандидата не требуются.

\subsection{Решение малой задачи}

Пусть
\[
    B_k=\widehat U_k\Theta_k\widehat V_k^\top,
    \qquad
    \Theta_k=\diag(\theta_1,\ldots,\theta_k).
\]
Тогда
\begin{equation}
\label{eq:projected-solution}
    y_{\lambda,k}
    =
    \widehat V_k
    \diag\left(
        \frac{\theta_i}{\theta_i^2+\lambda}
    \right)
    \widehat U_k^\top \rho e_1.
\end{equation}
SVD выполняется только для матрицы размера \((k+1)\times k\).
Полная SVD матрицы \(A\) не используется.

\section{Выбор \(\lambda\) по projected GCV}

Определим сглаживающую матрицу проецированной задачи
\[
    H_{k,\lambda}
    =
    B_k(B_k^\top B_k+\lambda I_k)^{-1}B_k^\top.
\]
Projected GCV задаётся функцией
\begin{equation}
\label{eq:pgcv}
    G_k(\lambda)
    =
    \frac{
        \norm{(I_{k+1}-H_{k,\lambda})\rho e_1}^{2}
    }{
        \left[
            \tr(I_{k+1}-H_{k,\lambda})
        \right]^2
    }.
\end{equation}
Множитель, не зависящий от \(\lambda\), можно опустить.

След вычисляется по сингулярным значениям \(B_k\):
\begin{equation}
\label{eq:dof}
    \tr(H_{k,\lambda})
    =
    \sum_{i=1}^{k}
    \frac{\theta_i^2}{\theta_i^2+\lambda}.
\end{equation}
Следовательно, вычисление \(G_k(\lambda)\) не требует:

\begin{itemize}[nosep]
    \item формирования \(A^\top A\);
    \item SVD полной матрицы \(A\);
    \item случайных trace-probe векторов;
    \item новых операций \(Av\) и \(A^\top u\) для каждого \(\lambda\).
\end{itemize}

Выбранный параметр:
\begin{equation}
\label{eq:lambda-choice}
    \lambda_k
    =
    \argmin_{\lambda>0}G_k(\lambda).
\end{equation}

Projected GCV оценивает prediction-риск проецированной линейной задачи.
Он не минимизирует напрямую ошибку направления
\[
    1-\left|
        \ip{\widehat\beta_\lambda/\norm{\widehat\beta_\lambda}}{\beta^\ast}
    \right|.
\]
Поэтому в синтетических экспериментах выбор необходимо проверять по
\(\lvert\cos(\widehat\beta_\lambda,\beta^\ast)\rvert\).
На реальных данных применяется cross-fitted ошибка локальных моментных уравнений.

\subsection{Weighted GCV}

Если обычный projected GCV систематически выбирает слишком большое \(\lambda\),
можно использовать
\begin{equation}
\label{eq:wgcv}
    G_{k,\omega}(\lambda)
    =
    \frac{
        \norm{(I_{k+1}-H_{k,\lambda})\rho e_1}^{2}
    }{
        \left[
            (k+1)
            -
            \omega\tr(H_{k,\lambda})
        \right]^2
    },
\end{equation}
где \(\omega>0\) калибруется экспериментально.
Weighted GCV не следует включать до обнаружения воспроизводимого смещения
обычного projected GCV.

После обнаружения систематического выбора слишком большого \(\lambda\)
используется вес
\begin{equation}
\label{eq:wgcv-full-weight}
    \omega_k=\frac{k+1}{JP}.
\end{equation}
Тогда знаменатель \eqref{eq:wgcv} равен
\[
    \frac{k+1}{JP}
    \left[
        JP-\tr(H_{k,\lambda})
    \right].
\]
Множитель \((k+1)/(JP)\) не зависит от \(\lambda\), поэтому минимум можно
эквивалентно искать для критерия
\begin{equation}
\label{eq:wgcv-full}
    G_{k,\mathrm{full}}(\lambda)
    =
    \frac{
        \norm{(I_{k+1}-H_{k,\lambda})\rho e_1}^{2}
    }{
        \left[
            JP-\tr(H_{k,\lambda})
        \right]^2
    }.
\end{equation}

\subsection{Cross-fitted проверка обновления}

Наблюдения один раз разбиваются на две непересекающиеся части.
Для каждой упорядоченной пары training--validation локальные коэффициенты
оцениваются только по training-моментам:
\[
    \widehat c_j^{(a)}(\beta)
    =
    \frac{
        \sum_p I_{jp}^{(a)}U_{jp}^{(a)}\beta
    }{
        \sum_p (U_{jp}^{(a)}\beta)^2+\eta
    }.
\]
Направление оценивается на противоположной части:
\begin{equation}
\label{eq:cross-fitted-moment-loss}
    L_{\mathrm{cf}}(\beta)
    =
    \frac{1}{N_{\mathrm{mom}}}
    \sum_{(a,b)\in\{(1,2),(2,1)\}}
    \sum_{j,p}
    \left(
        I_{jp}^{(b)}
        -
        \widehat c_j^{(a)}(\beta)U_{jp}^{(b)}\beta
    \right)^2.
\end{equation}
Кандидат применяется только если
\[
    L_{\mathrm{cf}}(\beta_{\mathrm{new}})
    \leq
    L_{\mathrm{cf}}(\beta_0).
\]
Для локальной инициализации первая пара систем локальных моментов сохраняется
как фиксированная валидационная цель. При случайной инициализации ранний
запрет по валидационному критерию не применяется: на разреженных локальных
моментах он может блокировать выход из плохого начального направления.

\section{Поиск минимума}

Унимодальность \(G_k(\lambda)\) в общем случае не гарантируется.
Поэтому используется следующая схема.

\begin{enumerate}[label=\arabic*.]
    \item Оценить \(\theta_1=\norm{B_k}\).
    \item Построить логарифмическую сетку
    \begin{equation}
    \label{eq:lambda-range}
        \lambda
        \in
        \left[
            10^{-8}\theta_1^2,\,
            10^{2}\theta_1^2
        \right].
    \end{equation}
    \item Вычислить \(G_k(\lambda)\) в 17--25 точках.
    \item Если минимум находится на границе, расширить диапазон на одну декаду.
    \item Для лучшего внутреннего интервала применить метод Брента к
    \(t=\log\lambda\).
    \item При нескольких сравнимых минимумах проверить два лучших значения
    по наблюдаемому validation-критерию.
\end{enumerate}

На следующей итерации используется continuation:
\[
    \lambda
    \in
    \left[
        \lambda_{\mathrm{prev}}/100,\,
        100\lambda_{\mathrm{prev}}
    \right],
\]
объединённый со спектральным диапазоном \eqref{eq:lambda-range}.
Под следующей итерацией здесь понимается следующий внутренний или внешний
запуск после изменения коэффициентов \(c_j\), весов либо \(U_j\).
Внутри одного процесса Голуба--Кахана значение
\(\lambda_{\mathrm{prev}}\) фиксировано для всех checkpoints и обновляется
только после возврата принятого кандидата.

\section{Нормировка направления}

После решения проецированной задачи вычисляется
\begin{equation}
\label{eq:normalize}
    \beta_{\mathrm{new}}
    =
    \frac{
        \widetilde\beta_{\lambda_k,k}
    }{
        \norm{\widetilde\beta_{\lambda_k,k}}
    }.
\end{equation}
Знак согласуется с предыдущим направлением:
\[
    \beta_{\mathrm{new}}
    \leftarrow
    \sign\!\left(
        \ip{\beta_{\mathrm{new}}}{\beta_0}
    \right)\beta_{\mathrm{new}}.
\]
Для контроля сходимости используется расстояние между прямыми
\begin{equation}
\label{eq:projective-distance}
    d_{\pm}(\beta_1,\beta_2)
    =
    \sqrt{
        2\left(
            1-\left|\ip{\beta_1}{\beta_2}\right|
        \right)
    }.
\end{equation}

Нормировка означает, что итоговый вектор не является точным минимумом
неограниченной евклидовой задачи \eqref{eq:ridge}. В single-index модели это
допустимо, поскольку оценивается направление. Параметр \(\lambda\) регулирует
величину и спектральный состав шага до проекции на сферу.

\section{Критерии остановки Krylov-процесса}

Проверки выполняются каждые \(\Delta k\) шагов после минимальной длины проекции.
Процесс останавливается, если два раза подряд выполнены условия
\begin{align}
\label{eq:stop-lambda}
    \left|
        \log_{10}\lambda_k
        -
        \log_{10}\lambda_{k-\Delta k}
    \right|
    &<0.1,
    \\
\label{eq:stop-direction}
    d_{\pm}
    \left(
        \beta_{\lambda_k,k},
        \beta_{\lambda_{k-\Delta k},k-\Delta k}
    \right)
    &<\varepsilon_{\mathrm{dir}},
    \\
\label{eq:stop-gcv}
    \frac{
        |G_k(\lambda_k)-G_{k-\Delta k}(\lambda_{k-\Delta k})|
    }{
        \max\{1,|G_k(\lambda_k)|\}
    }
    &<\varepsilon_{\mathrm{GCV}}.
\end{align}
Дополнительно выбранное \(\lambda_k\) должно лежать внутри диапазона поиска.

Практические начальные значения:
\[
    \Delta k=5,
    \qquad
    \varepsilon_{\mathrm{dir}}=10^{-3},
    \qquad
    \varepsilon_{\mathrm{GCV}}=10^{-2},
    \qquad
    k_{\max}=\min\{d,100\}.
\]
Значения \(k_{\max}\in\{50,100,200\}\) следует сравнить экспериментально.

\section{Интеграция в чередующуюся оптимизацию}

\begin{enumerate}[label=\textbf{Шаг \arabic*.}]
    \item Зафиксировать текущие локальные статистики \(I_j,U_j\) и направление
    \(\beta\).

    \item Обновить локальные коэффициенты, например
    \[
        c_j
        =
        \frac{
            \ip{I_j}{U_j\beta}
        }{
            \norm{U_j\beta}^{2}+\eta
        },
    \]
    где \(\eta>0\) стабилизирует почти вырожденные локальные задачи.

    \item Создать matrix-free операции \eqref{eq:matvec}--\eqref{eq:rmatvec}.

    \item Положить
    \[
        \beta_0=\beta,
        \qquad
        r_0=b-A\beta_0.
    \]

    \item Выполнить билинеаризацию Голуба--Кахана, одновременно контролируя
    стабилизацию \(\lambda_k\), projected GCV и направления.

    \item Выбрать \(\lambda_k\) по \eqref{eq:lambda-choice}, решить
    \eqref{eq:projected} и вычислить \(\beta_{\mathrm{new}}\) по
    \eqref{eq:normalize}.

    \item Если
    \[
        d_{\pm}(\beta_{\mathrm{new}},\beta)<\varepsilon_{\mathrm{inner}},
    \]
    завершить внутреннее чередование. Иначе положить
    \(\beta\leftarrow\beta_{\mathrm{new}}\) и вернуться к шагу 2.

    \item После изменения весов, направлений или \(U_j\) начать новый
    Krylov-процесс с warm start \(\beta_0=\beta_{\mathrm{previous}}\).
\end{enumerate}

После обновления \(c_j\) матрица \(A\) меняется. Старый Krylov-базис нельзя
считать базисом новой задачи. Warm start переносит \(\beta_0\) и
\(\lambda_{\mathrm{prev}}\), но не сами векторы \(V_k\), если не реализован
отдельный recycling-алгоритм.

\section{Обработка неудачных случаев}

\begin{center}
\begin{tabular}{>{\raggedright\arraybackslash}p{0.34\textwidth}
                >{\raggedright\arraybackslash}p{0.56\textwidth}}
\toprule
Ситуация & Действие\\
\midrule
Минимум GCV на границе &
Расширить диапазон на одну декаду и повторить поиск.\\

\(\lambda_k\) нестабилен при росте \(k\) &
Продолжить Голуба--Кахана до стабилизации или \(k_{\max}\).\\

Несколько локальных минимумов &
Проверить два лучших значения по cross-fitted moment loss.\\

\(\lambda\) слишком велико, \(\beta_{\mathrm{new}}\approx\beta_0\) &
Сравнить \(\lambda/10,\lambda,\lambda\cdot10\); уменьшить параметр при
улучшении validation-критерия.\\

\(\lambda\) слишком мало и решение чувствительно к \(k\) &
Увеличить нижнюю границу либо ограничить вклад малых Ritz-значений.\\

Нулевой или почти нулевой \(r_0\) &
Сохранить \(\beta_0\); новый Krylov-процесс не требуется.\\

NaN, Inf или вырожденная малая задача &
Использовать \(\lambda_{\mathrm{prev}}\), увеличить \(\eta\) и проверить
масштабирование \(I_j,U_j\).\\

Известна дисперсия шума &
Заменить projected GCV на projected UPRE или discrepancy principle.\\
\bottomrule
\end{tabular}
\end{center}

\section{Стоимость и память}

Пусть \(T_{\mathrm{op}}\) --- стоимость одного вызова \(Av\) или \(A^\top u\).

Один hybrid-запуск с \(k\) шагами требует
\[
    O(kT_{\mathrm{op}})
\]
для операторных вычислений. После построения проекции обработка \(L\) кандидатов
требует
\[
    O(k^3+Lk)
\]
операций: \(O(k^3)\) на SVD \(B_k\) и \(O(Lk)\) на значения критерия.

При отдельных LSMR-запусках по сетке стоимость составляет
\[
    O(LkT_{\mathrm{op}}).
\]
Поэтому hybrid-схема устраняет множитель \(L\) перед дорогими
matrix-free операциями.

Дополнительная память:
\[
    O(dk+k^2)
\]
для \(V_k\) и \(B_k\).
Хранить \(U_{k+1}\) полностью не обязательно, если не требуется повторная
ортогонализация слева. При \(d=1000\) и \(k=100\) матрица \(V_k\) занимает
около \(0.8\) МБ в формате \texttt{float64}.

\section{Рекомендуемая конфигурация}

Основная схема:
\[
\boxed{
\begin{gathered}
    \delta=\beta-\beta_0,\qquad r_0=b-A\beta_0
    \\
    \text{Golub--Kahan}
    \longrightarrow
    B_k,V_k
    \longrightarrow
    \text{projected ridge}
    \\
    \longrightarrow
    \text{projected GCV}
    \longrightarrow
    \lambda_k
    \longrightarrow
    \beta_{\mathrm{new}}
    =
    \operatorname{normalize}
    \left(
        \beta_0+V_ky_{\lambda_k,k}
    \right).
\end{gathered}
}
\]

Projected GCV используется как дешёвый основной селектор.
В синтетических тестах его необходимо сравнивать с параметром
\[
    \lambda_{\mathrm{dir}}
    =
    \argmax_\lambda
    \left|
        \cos(\widehat\beta_\lambda,\beta^\ast)
    \right|.
\]
Если GCV систематически ухудшает качество направления, окончательный выбор
следует выполнять по cross-fitted moment validation, сохраняя один
Krylov-процесс для дешёвого вычисления кандидатов.

\section{Замечание о реализации}

Готовая функция \texttt{scipy.sparse.linalg.lsmr} не возвращает базис \(V_k\)
и бидиагональную матрицу \(B_k\). Для hybrid-метода требуется собственная
реализация билинеаризации Голуба--Кахана либо библиотека, предоставляющая
доступ к проецированной задаче. Matrix-free функции \texttt{matvec} и
\texttt{rmatvec} из текущего LSMR-солвера можно использовать без изменения.

\begin{thebibliography}{9}

\bibitem{PaigeSaunders1982}
C.~C. Paige, M.~A. Saunders.
\newblock LSQR: An algorithm for sparse linear equations and sparse least squares.
\newblock \emph{ACM Transactions on Mathematical Software},
8(1):43--71, 1982.

\bibitem{FongSaunders2011}
D.~C.-L. Fong, M.~A. Saunders.
\newblock LSMR: An iterative algorithm for sparse least-squares problems.
\newblock \emph{SIAM Journal on Scientific Computing},
33(5):2950--2971, 2011.

\bibitem{GolubHeathWahba1979}
G.~H. Golub, M.~Heath, G.~Wahba.
\newblock Generalized cross-validation as a method for choosing a good ridge parameter.
\newblock \emph{Technometrics},
21(2):215--223, 1979.

\bibitem{OLearySimmons1981}
D.~P. O'Leary, J.~A. Simmons.
\newblock A bidiagonalization-regularization procedure for large scale
discretizations of ill-posed problems.
\newblock \emph{SIAM Journal on Scientific and Statistical Computing},
2(4):474--489, 1981.

\bibitem{ChungNagyOLeary2008}
J.~Chung, J.~G. Nagy, D.~P. O'Leary.
\newblock A weighted-GCV method for Lanczos-hybrid regularization.
\newblock \emph{Electronic Transactions on Numerical Analysis},
28:149--167, 2008.

\bibitem{Hansen1992}
P.~C. Hansen.
\newblock Analysis of discrete ill-posed problems by means of the L-curve.
\newblock \emph{SIAM Review},
34(4):561--580, 1992.

\end{thebibliography}

\end{document}
